\subsection{P-19. Balanced two-pair evaluation hits every multilinear width-three cubic}
\label{P-19}

\paragraph{Statement.}
Fix pairwise distinct nodes $a_1,\ldots,a_n$ over a field $E$.
Every nonzero homogeneous multilinear cubic having ROABP width
at most three in some variable order is hit by the balanced
two-pair map
\[
 R_i=t+\frac{z_1}{a_i-y_1}+\frac{z_2}{a_i-y_2}.
 \tag{P19.1}
\]
The variable order need not be known. The statement holds in
every characteristic and for every number of variables.
Consequently the denominator-safe finite grid of I-40 gives
a uniform hitting set of at most
\[
 20(n+1)^2
 \tag{P19.2}
\]
points over a working field of size at least
$\max\{2n+1,4\}$, with the extension-input convention stated
there. By multilinear reversal, the corresponding formal
nonvanishing assertion also holds for homogeneous degree
$n-3$.

\paragraph{Inputs.}
We use the prime balanced one-pair ideal $\mathcal I_1$, its
absence of nonzero linear forms, and pole extraction
$\partial_i\mathcal I_2\subseteq\mathcal I_1$ for multilinear
cubics, as specified in P-17--P-18. We also use the proved
minimal-counterexample interval reduction in P-18, part~2,
and the established cubic exterior characterization of
Hu, van Melkebeek, and Morgan (2024), Definition~9.5, p.~55,
and Theorem~9.13, pp.~61--62 (P-B20). No new characterization
of the full vanishing ideal is assumed.

\begin{proof}
Suppose a counterexample exists. Among counterexamples on
subsets of its variables, choose one with the fewest essential
variables, remove unused variables, and relabel a witnessing
width-three order as $1,\ldots,N$. The corresponding nodes
remain distinct. Denote this nonzero cubic by $f$.

P-18, part~2, supplies consecutive position intervals
\[
 A=[1,\alpha],\qquad B=[\beta,\gamma],\qquad
 C=[\delta,N],\qquad
 \alpha\le\beta\le\gamma\le\delta,
 \tag{P19.3}
\]
such that every support triple has its first, second, and third
positions in $A,B,C$, respectively. There are no inessential
positions. The intersections $A\cap B$ and $B\cap C$ have
at most one position each. This use of P-18 retains its
minimal-counterexample hypothesis; no arbitrary coordinate
specialization is asserted to preserve the balanced ideal.

We first show that at most one position of $B$ lies outside
$A\cup C$. Suppose instead that there are two such positions
$j<k$. Each is essential and can occur only as the second
position of a support triple. Therefore
\[
 Q_j=\partial_j f
   =\sum_{i\in A,\,\ell\in C}q^{(j)}_{i\ell}X_iX_\ell,
 \qquad
 Q_k=\partial_k f
   =\sum_{i\in A,\,\ell\in C}q^{(k)}_{i\ell}X_iX_\ell
 \tag{P19.4}
\]
are nonzero bilinear quadratics. Here every $i\in A$ precedes
$j$, and every $\ell\in C$ follows $k$, so the two extreme
variable sets are disjoint.

Both quadratics belong to $\mathcal I_1$ by pole extraction.
Their displayed bilinear coefficient matrices each have rank
at least two. Indeed, rank one would express the corresponding
nonzero quadratic as a product of two nonzero linear forms.
Primality of $\mathcal I_1$ would put one of those forms in
$\mathcal I_1$, contrary to the absence of linear equations.
This argument applies over $E$ and uses no division by the
characteristic.

Take the coefficient cut immediately after position $j$.
In its left-degree-two block, restrict rows to $X_iX_j$ with
$i\in A$ and columns to $X_\ell$ with $\ell\in C$.
The resulting submatrix is the coefficient matrix of $Q_j$,
so this degree block has rank at least two. In its
left-degree-one block, restrict rows to $X_i$ with $i\in A$
and columns to $X_kX_\ell$ with $\ell\in C$. This gives the
coefficient matrix of $Q_k$, so this block also has rank at
least two. Since $f$ is homogeneous, these degree blocks
have disjoint row and column index sets. The full coefficient
cut therefore has rank at least four, contradicting the
witnessing width bound.

We have proved
\[
 |B\setminus(A\cup C)|\le1.
 \tag{P19.5}
\]
Together with the two possible singleton intersections in
(P19.3), this implies $|B|\le3$.

Now put
\[
 \mathscr K=\ker
 \begin{pmatrix}
 1&\cdots&1\\ a_1&\cdots&a_N\\ a_1^2&\cdots&a_N^2
 \end{pmatrix},
 \qquad
 W=\sum_{|S|=3}\frac{c_S}{\Delta_S}e_S
       \in\bigwedge^3\mathscr K,
 \tag{P19.6}
\]
where $c_S$ is the coefficient of $X_S$ in $f$ and
$\Delta_S$ is the consistently ordered three-node Vandermonde
determinant. The cited characterization gives the membership
in (P19.6), and $W\ne0$ because all $\Delta_S$ are nonzero.

Let $\pi:E^N\to E^{[N]\setminus B}$ be coordinate projection.
Its restriction to $\mathscr K$ is injective. A vector in its
kernel would be supported on at most the three coordinates
of $B$, whereas any at most three columns of the displayed
Vandermonde matrix are independent. The exterior cube of an
injective linear map is injective as well. But every support
triple of $f$, and hence of $W$, meets $B$ by (P19.3).
It follows that
\[
 (\bigwedge^3\pi)(W)=0,
 \tag{P19.7}
\]
contradicting this injectivity and $W\ne0$.

Thus no cubic counterexample exists. The interpolation
argument of I-40 at total degree three gives (P19.2), with
the stated selector and amplitude field sizes. Multilinear
reversal preserves both every coefficient-cut rank and the
balanced rational family, as established in P-17, so it gives
the degree-$n-3$ assertion. This completes the proof.
\end{proof}

\paragraph{Scope and remaining higher-degree problem.}
This proves the full cubic statement, including both shared
boundary cases left by P-18. It does not prove the full
multilinear two-pair theorem in arbitrarily many variables,
nor does it identify balanced two-pair evaluation with an
unweighted ordinary centered SV map. For a remaining
multilinear width-three annihilator, the lowest nonzero
homogeneous degree must be at least four, and the highest
must be at most $n-4$. For the first claim use the known
inheritance of width by the lowest homogeneous component
(HMM Proposition~8.3, p.~42), together with the absence of
degrees below three and the theorem just proved; apply the
same argument after multilinear reversal for the second.
The higher-degree and growing-width sampling gates remain
separate. No historical priority claim is made.
